\section{EXPERIMENTAL DESIGN}\label{sec:method}

\paragraph{Testbed}
A single node runs the full stack in containers: PostgreSQL, Redpanda, Airflow, and OPA.
Table~\ref{tab:testbed} lists the hardware, software versions, and fixed parameters. The policy
engine is pinned to a specific release, and Section~\ref{sec:results} shows why that pin matters.
Runs execute strictly sequentially, since the resource-contention scenario applies real CPU stress
and concurrent runs would fold host load into the measurement.

\begin{table*}[t]
\centering
\caption{Experimental testbed and fixed design parameters.}
\label{tab:testbed}
\small
\begin{tabular}{@{}lp{0.74\linewidth}@{}}
\toprule
Item & Setting \\
\midrule
Host & One Apple M2 machine (arm64), 8 cores, 8\,GB memory, macOS 15.7.2; runs are strictly sequential \\
Software & Docker 29.6.1; PostgreSQL 16.6; Redpanda v24.2.18; Apache Airflow 2.10.5 (Python 3.11); OPA 0.68.0 \\
Batch data & Downscaled TPC-DS-derived synthetic fact table, 20{,}000 rows \\
Live model & Reasoning agents: \texttt{nvidia/nemotron-3-ultra-550b-a55b}; monitoring: \texttt{nvidia/nemotron-3-super-120b-a12b}; hosted OpenAI-compatible endpoint, temperature 0 \\
LLM budget & \NBudgetCalls{} calls and \NBudgetTokens{} tokens per run \\
Timing & 120\,s warm-up, 300\,s control loop, 5\,s settle; control tick 15\,s \\
Human model & Log-normal delay, median \NHumanMedian{}\,s, $\sigma=\NHumanSigma{}$ \\
Cost model & Static allocation 8 units, right-sized 3 units, horizon 300\,s; 0.05 per compute unit-second, 0.01 per GB-hour \\
Baseline latencies & \texttt{rule\_based} resolves covered faults 30\,s after detection; \texttt{autoscale} 20\,s after detection \\
Policy limits & Rate limit of 5 actions per 10 minutes; default cost budget of 100 units per run \\
\bottomrule
\end{tabular}
\end{table*}

\paragraph{Scenarios and configurations}
Four injected faults cross eight configurations. The faults, schema drift (a column's type changes),
upstream delay, resource contention, and an ingress burst, pair with \texttt{baseline} (static
orchestration plus the simulated human), \texttt{rule\_based} (deterministic automation),
\texttt{autoscale} (threshold autoscaling), \texttt{full} (all four agents), and four ablations.
Three ablations add exactly one agent on top of monitoring, namely \texttt{recovery\_only},
\texttt{optimization\_only}, and \texttt{schema\_only}, while the fourth, \texttt{monitor\_only},
runs monitoring on its own. None of the three non-agent configurations ever calls a model.

\paragraph{Non-agent baselines}
The three non-agent configurations model behavior rather than run a deployed controller. In
\texttt{baseline}, every fault escalates to the simulated on-call. In \texttt{rule\_based}, a fault
inside the rule set, upstream delay, resource contention, or ingress burst, resolves 30\,s after
detection, while schema drift, which demands reasoning, escalates to the human. In
\texttt{autoscale}, the reactive scaler resolves resource contention and ingress bursts 20\,s after
detection, and every data fault escalates. Their recovery times therefore encode two assumed
remediation latencies rather than measuring an actual rule engine or autoscaler. They are idealized
comparators: speed depends on those latencies, while any fault type outside their coverage still
falls to the human, which bounds what they can achieve.

\paragraph{Matrix and arms}
Each (configuration, scenario) cell carries $\NNBaseline{}$ replicates for the baselines and
\texttt{full}, fewer for the ablations; Table~\ref{tab:main} gives the realized $n$ per cell. Because
only agent configurations call a model, the campaign splits into three independently resumable arms
with separate result directories, guaranteeing that live and mock conditions can never be silently
pooled. Arm A holds the live-model agent configurations; arm B holds the non-agent configurations
with the model disabled, whose results are identical live or mock by construction; arm C runs
\texttt{full} under a deterministic mock as the direct counterpart to arm A's \texttt{full}. The
three arms contain \NRunsArmA{}, \NRunsArmB{}, and \NRunsArmC{} runs, \NRunsTotal{} in total, and
every run has a 120\,s warm-up, a 300\,s control loop after fault injection, and a 5\,s settle.

\paragraph{Live model and accounting}
The live arm calls a remote language model at temperature zero through an OpenAI-compatible
endpoint, with the reasoning agents sharing one model and monitoring using a smaller, faster one.
Model identifiers are recorded per run in the run manifest, and a per-run budget caps model calls
and tokens.

Two properties of the accounting matter for validity. First, a snapshot-identical proposal comes
from an in-run cache, so summing per-action token columns would count one real call several times
over; we instead count real calls, cache hits, degraded proposals, invalid outputs, tokens, and
latency only at the moment a call is or is not actually made, and a proposal counts as degraded
whenever the budget is exhausted, the provider is unavailable, or an earlier degradation is
replayed. Second, a transient provider failure is cached against the snapshot that triggered it, so
a single rate-limit episode can keep degrading every later identical snapshot unless the degraded
count is tracked separately, exactly what happened during the project's own pilot
(Section~\ref{sec:lessons}) and why the accounting takes this form.

\paragraph{Metrics}
We report the five metrics defined in Section~\ref{sec:problem}: mean time to recovery (MTTR,
Eq.~\ref{eq:mttr}), operational cost under a disclosed model (Eq.~\ref{eq:cost}), manual
interventions, data freshness, and \emph{decision correctness} (Eq.~\ref{eq:decision}). Manual
interventions are a count of rows in the telemetry table every configuration writes to, and
freshness is a lag in seconds read from the same telemetry.

\paragraph{Seeds and provenance}
Every run's seed is $\mathrm{sha256}(\texttt{config:scenario:replicate}) \bmod 2^{32}$. Hashing the
configuration name in means runs matched on (scenario, replicate) share no fault or human-latency
randomness across configurations, so the matching is nominal rather than paired; we treat the
samples as independent and use the rank-sum test as primary, keeping the signed-rank test on the
matched index only for continuity with earlier analysis. Each manifest line records the profile, the
model mode and identifiers (never credentials), timings, the human-latency and provisioning
parameters, budgets, the code revision, and whether the experiment code was uncommitted. Analysis
refuses to merge runs whose provenance differs unless given an explicit, documented override.

\paragraph{Statistics}
We report medians with interquartile ranges, plus bootstrap 95\% intervals for medians and for
differences between them. Cliff's $\delta$ always comes with its own bootstrap interval, read
against the conventional thresholds: $|\delta|<0.147$ negligible, $0.147$--$0.33$ small,
$0.33$--$0.474$ medium, and $\geq 0.474$ large~\cite{cliff1993}. Mann--Whitney $U$ tests carry a
Holm correction across the full family of comparisons, and proportions carry Wilson intervals. Every
comparison in Section~\ref{sec:results} is exploratory: the campaign was designed and run before any
single comparison was drafted, not as a pre-registered confirmatory study.

\paragraph{Adversarial evaluation}
A generated corpus probes the policy layer directly, with an exhaustive grid over every legal
(agent, action) pair and the decision-relevant context values, including boundary cases such as 4,
5, and 6 recent actions and a cost equal to or infinitesimally above budget; attacker-controlled
parameter extremes priced by the real gate; hostile strings in free-text fields; illegal proposals
sent straight to the policy engine, bypassing the contract; and a seeded fuzz. Expected verdicts
come from an independently written Python statement of the \emph{documented} policy semantics,
sharing no code with the Rego under test; the independence is one of implementation, not authorship.
The threat model gives the agent control only over its own proposal, leaving the surrounding policy
context system-built and out of reach.

\paragraph{Sensitivity}
Sampled human latency follows $\mathrm{median}\cdot e^{\sigma Z}$ with $Z$ fixed by the seed, a
scale family under which any all-human outcome scales exactly linearly in the assumed median; we
confirm the linearity directly against the sampler and use it to derive, in closed form, the human
median at which each configuration stops beating the baseline. For the cost model, we sweep the
provisioning parameters and report where the break-even falls.
